% ch10_d §10.8–10.9 (书页 353–361 / p367–p375)
%--C-TAIL--
%JOIN%
我们把均匀磁化率 (10.42) 推广为一个依赖于 $q$ 的磁响应函数。沿着 §5.1 的思路，在导致磁化率 $\chi_0(\mathbf{q})$（即式 (10.51)）的那类计算中，纳入一个形式如 (10.38) 的交换相互作用是很容易的。所得结果实际上与式 (10.42) 明显相似，即

\begin{equation*}
\chi(\mathbf{q}) = \frac{\chi_0(\mathbf{q})}{1 - \tfrac{1}{2} U (g\beta s)^{-2} \chi_0(\mathbf{q})}.
\tag{10.67}
\end{equation*}

现在，如果这个表达式的分母为零，金属的假定非磁态便是不稳定的。磁有序的 Overhauser 条件是

\begin{equation*}
\tfrac{1}{2} U (g\beta s)^{-2} \chi_0(\mathbf{q}) > 1
\tag{10.68}
\end{equation*}

其中铁磁性的条件 (10.43) 正是 $q = 0$ 的特殊情形。

通常的自由电子间交换相互作用能否在普通金属中导致稳定的自旋密度波（spin density waves），这个问题尚未得到确定的解决。然而在过渡金属中，已知 $U$ 相当大，而待代入 $\chi_0(\mathbf{q})$ 的公式 (10.51) 的电子能态 $\mathcal{E}(\mathbf{k})$ 由能带结构主导。例如，由 §5.4 显然可知，一个横跨费米面相互平行的平坦区域的波矢 $\mathbf{q}$ 可能诱导出很大的 $\chi_0(\mathbf{q})$ 值，从而使 (10.68) 得以满足。在这种情形下，反铁磁有序可能异常复杂：自旋密度的变化等价于锥面上自旋矩的螺旋形排布，其重复距离与底层晶格无关。

\section{Ising 模型}

在讨论铁磁性和反铁磁性时，我们对 Heisenberg 哈密顿量 (10.29) 的态采用了非常粗糙的近似。实际上，我们假定除正在处理的那个自旋之外，每个自旋都取其平均值；我们忽略了相邻格点上自旋之间的涨落与关联。此外，我们还忽略了 $\mathsf{S}_l$ 是算符这一事实——它们不能换成各自的平均期望值。

自旋哈密顿量的本征值分布实际上非常复杂，试图对系统的热力学给出精确的讨论，是固体理论的主要课题之一。这一理论沿两个方向发展。我们可以假定系统接近其基态（铁磁的或反铁磁的），然后讨论低能激发。这就是自旋波（spin waves）理论，将在 §10.11 中概述。

另一类近似是略去自旋算符的非对角元，只考虑它们沿某个固定方向——通常是 $H$ 的方向——的分量。换言之，我们假定系统的能级由下式给出

\begin{equation*}
\mathcal{E} = -\sum_{ll'} J_{ll'} \sigma_l \sigma_{l'} - \beta H \sum_l \sigma_l,
\tag{10.69}
\end{equation*}

%FIG{185}: Ising 模型的一个组态 (a) 可以代表：(b) 自旋的一种排布；(c) 二元合金中原子的一种排布；(d) "格气"的一个组态。

其中我们在每个格点 $l$ 上定义一个只能取 $+1$ 或 $-1$ 的“自旋”量子数 $\sigma_l$。在多数情况下我们还假定 $J_{ll'}$ 只在格子中最近邻之间起作用。

这就叫作 Ising 模型。它不可能是铁磁或反铁磁系统的精确描述，但由于它绕开了求自旋集合的精确本征函数这一极为困难的问题，使人们得以集中处理这类系统的态的统计分布这一组合问题。

Ising 模型确实更精确地描述了另一类系统——合金：两种元素 A 和 B 可以在一组固定的格点位置上相互替换。我们可以说 $\sigma_l = +1$ 对应于第 $l$ 个格点上的一个 A 原子，而 $\sigma_l = -1$ 规定该格点上是一个 B 原子。“交换积分” $J$ 这时应解释为

\begin{equation*}
J = \tfrac{1}{2}\bigl\{ \mathcal{V}_{AB} - \tfrac{1}{2}(\mathcal{V}_{AA} + \mathcal{V}_{BB}) \bigr\},
\tag{10.70}
\end{equation*}

其中 $\mathcal{V}_{AB}$ 是相邻格点上 A 原子与 B 原子之间的相互作用能，减去的则是一个 AA 对或一个 BB 对的平均能量。

这类合金系统可以作实验研究，而且发现它们的确表现出一个过渡温度，类似于居里温度或 Néel 温度。在转变点以上，合金是无序的；温度较低时它趋于有序，或者是通过分凝出近乎纯 A 与近乎纯 B 的材料（这对应于 $J > 0$ 的情形，即铁磁性），或者是通过形成 A、B 原子规则地分布在相互穿插的子晶格上的区域——这是反铁磁性的类似物。有序–无序转变（order–disorder transition）与铁磁或反铁磁转变的差别在于：前一情形中 A 类与 B 类原子的总数是固定的，而“向上”与“向下”的自旋却可以自由地相互转变。这一差别在理论上可以形式地加以顾及，因此许多结果对两类系统普遍成立。

Ising 模型可以应用的另一个问题是“格气”（lattice gas）或“格液体”（lattice liquid）问题：把一些能在短距离上相互作用的原子放到一组格点上，并留下空格点。这在形式上与合金问题等价。然而，它是否真正贴切地适用于人们有时让它去描述的液–气转变的热力学，却令人怀疑。

不过，Ising 模型的重要性并不在于描述个别的物理效应；它是一个数学上易于处理的模型，描写一个应当表现出合作现象（co-operative phenomena）和相变的系统。在这个意义上，它更多地属于统计力学的一般理论，而不属于固体理论。然而，有某些原理对晶体格子的一般理论也很重要，我们现在就试图阐明它们。

\section{组合方法}

Ising 模型的优点是，其热力学性质的计算可以归结为一个组合问题。例如，假设我们要求配分函数（partition function）

\begin{equation*}
Z = \sum e^{-\mathcal{E}/kT},
\tag{10.71}
\end{equation*}

其中求和遍及所有“组态”，即遍及格子的 $N$ 个格点上 $+$ 自旋与 $-$ 自旋的各种排布。这个和看起来非常复杂，但形式上它只需涉及两个变量，例如 $N_A$——“向上”自旋[或 A 原子]的数目——与 $N_{AB}$——相邻反平行自旋[或 AB 对]的数目。

为看清这一点，我们把能量 (10.69) 用各种类型的键的总数表出：

\begin{equation*}
\mathcal{E} = J(N_{AB} - N_{AA} - N_{BB}) - \beta H (N_A - N_B).
\tag{10.72}
\end{equation*}

但这些变量并非相互独立。显然有

\begin{equation*}
N_A + N_B = N.
\end{equation*}

同样显然的是，止于 A 格点的键至多只有 $pN_A$ 条，这里 $p$ 是格子的配位数（co-ordination number），即给定格点的最近邻数目。每条这样的键需要在 $N_{AA}$ 中被计及两次、在 $N_{AB}$ 中被计及一次。于是 $2N_{AA} + N_{AB} = pN_A$；类似地 $2N_{BB} + N_{AB} = pN_B$。代入这些条件，我们发现 (10.72) 变成

\begin{equation*}
\mathcal{E} = -\tfrac{1}{2} pNJ + 2N_{AB} J - (2N_A - N)\beta H.
\tag{10.73}
\end{equation*}

现在设 $g(N; N_A, N_{AB})$ 表示把 $N_A$ 个“向上”自旋放到格子上、且恰有 $N_{AB}$ 对反平行近邻的全部放法数。所有这些组态都具有相同的能量；我们可以把配分函数 (10.71) 写成

\begin{align*}
Z &= y^{\frac{1}{2}N} z^{-\frac{1}{2}pN} \sum_{N_A, N_{AB}} g(N; N_A, N_{AB})\, y^{N_A} z^{N_{AB}} \\
&= y^{\frac{1}{2}N} z^{-\frac{1}{2}pN} \Lambda_N(y, z),
\tag{10.74}
\end{align*}

其中变量 $y$ 和 $z$ 与磁场及相互作用 $J$ 相联系：

\begin{equation*}
y \equiv e^{2\beta H/kT}; \quad z \equiv e^{-2J/kT}.
\tag{10.75}
\end{equation*}

系统的全部有趣性质现在都集中在函数 $\Lambda_N$ 之中；由它可以导出适用于“铁磁”与“反铁磁”系统、或合金相与正规溶液（regular solutions）理论的公式。

但是，计算 (10.74) 中组合因子的问题一般是不可解的；已知的只是近似公式。对给定的 $N_A$ 值，算出组态的总数并不困难。由初等代数，这必是把 $N_A$ 个物体放到 $N$ 个位置上的排布方式总数

\begin{equation*}
\sum_{N_{AB}} g(N; N_A, N_{AB}) = \frac{N!}{N_A! (N - N_A)!}.
\tag{10.76}
\end{equation*}

困难在于把这个总数按 $N_{AB}$ 的各个取值分开。

一个非常粗糙的近似，是把对 $N_{AB}$ 的单独求和换成假定它取其平均值。也就是说，我们把 $N_A$ 个 A 原子和 $N_B$ 个 B 原子无规地分布在格子的格点上，并计算 $\tfrac{1}{2}pN$ 条键中任一条一端为 A 原子、另一端为 B 原子的概率。结果当然是

\begin{align*}
\langle N_{AB} \rangle &= 2\, \frac{N_A}{N}\, \frac{N_B}{N}\, \tfrac{1}{2} pN \\
&= \frac{p N_A N_B}{N}.
\tag{10.77}
\end{align*}

把 (10.76) 和 (10.77) 代入 (10.74)，得

\begin{equation*}
\Lambda_N \approx \sum_{N_A} \frac{N!}{N_A! N_B!}\, y^{N_A} z^{p N_A N_B / N}.
\tag{10.78}
\end{equation*}

由此可以计算自由能：从对数中取出最大的项，并用 Stirling 公式作渐近展开。于是

\begin{align*}
F &= -kT \ln \Lambda_N \\
&\approx kT \Bigl\{ -N \ln N + N_A \ln N_A + N_B \ln N_B - N_A \ln y - p\, \frac{N_A N_B}{N} \ln z \Bigr\}.
\tag{10.79}
\end{align*}

%FIG{186}: 对 $N=5$、$N_A=3$ 的链作组态计数。此系统的 $\langle N_{AB}\rangle = 2.4$。

自由能取极小（也即求和中出现极大项）的条件，可通过对 $N_A$ 求导得到（当然要取 $N_B = N - N_A$）。利用 (10.75)，我们得到

\begin{equation*}
\ln N_A - \ln N_B - \frac{2\beta H}{kT} - \frac{(N_A - N_B)}{N}\, \frac{2pJ}{kT} = 0,
\end{equation*}

即

\begin{equation*}
\frac{N_A}{N_B} = \exp \Bigl\{ \frac{2\beta H}{kT} + \Bigl( \frac{N_A - N_B}{N} \Bigr) \frac{2pJ}{kT} \Bigr\},
\tag{10.80}
\end{equation*}

或

\begin{equation*}
\Bigl( \frac{N_A - N_B}{N} \Bigr) = \tanh \Bigl\{ \frac{\beta H}{kT} + \Bigl( \frac{N_A - N_B}{N} \Bigr) \frac{pJ}{kT} \Bigr\}.
\tag{10.81}
\end{equation*}

如果把 $(N_A - N_B)/N$ 认作 $\langle \mu \rangle$，即系统的平均净极化，并像 (10.33) 那样把 $pJ$ 表为一个内场，那么我们就回到了铁磁转变的居里–外斯条件 (10.26)。在合金理论中，这被称为 Bragg--Williams 近似（Bragg--Williams approximation）。

%FIG{187}: 准化学近似下的比热，格子的配位数为 p = 12。

上面这个初等内场公式的推导暴露了它的局限性；它对相邻自旋之间的关联只字未提。例如在铁磁体中，由于自旋对有平行排列的倾向，我们可以预期随机平均值 (10.77) 高估了实际平衡组态中 AB 对的数目。

改进这一点的一个有趣途径是 Guggenheim 的准化学近似（quasi-chemical approximation）。假设我们在无规混合的基础上算出与 (10.77) 类似的其他平均值，就得到

\begin{equation*}
\frac{\langle N_{AB} \rangle^2}{\langle N_{AA} \rangle \langle N_{BB} \rangle} = 4.
\tag{10.82}
\end{equation*}

但如果我们处理的是分子之间的化学平衡，如同化学反应

\begin{equation*}
AA + BB \rightleftharpoons 2AB,
\tag{10.83}
\end{equation*}

所表达的那样，那么我们知道，不同组元的平衡浓度应由下式给出

\begin{equation*}
\frac{\langle N_{AB} \rangle^2}{\langle N_{AA} \rangle \langle N_{BB} \rangle} = 4\, \frac{\bigl( e^{-\mathcal{E}_{AB}/kT} \bigr)^2}{e^{-\mathcal{E}_{AA}/kT}\, e^{-\mathcal{E}_{BB}/kT}},
\tag{10.84}
\end{equation*}

其中 $\mathcal{E}_{AB}$ 等等是 AB 分子等等的能量。用 (10.75) 的记号，这提示 (10.84) 的右端平均而言应等于 $4z^2$。对 (10.73) 稍作变换，便可算出 $\langle N_{AB} \rangle$，并用它代替 (10.77) 放入 (10.78)。

这个计算我们就不再往下做了，只想指出：在由下式定义的转变温度以下，系统表现出自发极化（铁磁性或反铁磁性）

\begin{equation*}
\frac{kT_c}{pJ} = \frac{2/p}{\ln(1 - 2/p)},
\tag{10.85}
\end{equation*}
（译注：按准化学近似的标准结果，右端应为 $-2/p\big/\ln\frac{p}{p-2}$，原书疑漏负号。）

在配位数 $p$ 很大的极限下，它应趋于 (10.28)（在我们现在的记号中 $k\theta = pJ$）。我们还可以近似地计算转变温度附近的比热；它表现出一个尖锐的奇点，在高温一侧有一间断，如图 187 所示。

碰巧，准化学近似等价于另一个看上去颇不相同的方法，它出自 Bethe 和 Peierls。这是把内场论证加以推广的一种尝试。取一个自旋集团——例如中心一个位于 $l = 0$ 的自旋，以及它的 $p$ 个近邻。从这个集团的能量中取 (10.69) 的下列各项：

\begin{equation*}
\mathcal{E}_p = -\sum_{j=1}^{p} J \sigma_0 \sigma_j - \beta H \sigma_0 - \beta H_I \sum_{j=1}^{p} \sigma_j.
\tag{10.86}
\end{equation*}

换言之，我们假定中心自旋感受到实际的外场 $H$，并且显式地与近邻的自旋相连；而对集团的外围格点，我们则假定可以把它们与环境的相互作用当作一个内场 $H_I$ 的效应来处理。

%FIG{188}: Bethe–Peierls 集团。

很容易为这个集团写出完整的配分函数，因为它只含 $p$ 项。于是可以算出集团中心原子的平均取向。但中心格点并不是格子中什么特权位置；$\langle \sigma_0 \rangle$ 必须与 $\langle \sigma_j \rangle$ 相同，后者是集团边缘上一个自旋的平均矩，它同样可以写成 $H_I$ 的函数。这就给出了内场 $H_I$ 的一个自洽条件：即使 $H = 0$，在低于由 (10.75) 定义的转变温度 $T_c$ 的温度下，$H_I$ 也可以不为零。

Bethe 方法作为一种物理上直观的程序是有吸引力的，甚至可以让它对更现实的模型给出结果：引入真正的自旋算符来代替 (10.86) 的“Ising 自旋变量”，并对这个集团哈密顿量的能级计算配分函数。或者，我们也可以在第一近邻的简单集团之外再引入更多的壳层，把问题解到更高的近似。

这些更高阶的近似，实际上是想更准确地确定 $\Lambda_N(y, z)$ 幂级数展开的系数。可以证明，每一个系数——无论是低温下有效的按 $z$ 的幂的展开，还是高温下有效的按

\begin{equation*}
w = \tanh \frac{J}{kT},
\tag{10.87}
\end{equation*}

的幂的展开——都依赖于数出在格子的“键”上能够描出的、具有给定边数的多边形，这颇像非理想气体理论中的 Mayer 集团和展开，或微扰论中的“图形”。这种计数原则上足够容易，实际做起来却极其繁难。

此外，这样的展开还不能完全使人相信：系统在某个确定温度的热力学性质中真的存在奇点；而 Bragg--Williams 方法和 Bethe 方法的封闭解析公式也填补不了这个空缺。这些方法\emph{假定}存在一个内场，或者假定系统有一个平均说来在每个格点上都相同的净极化。换言之，它们假定系统具有长程序（long-range order）——也就是说，只要知道一个格点上自旋的方向，我们就能猜出任何其他格点上它的方向，不管有多远。\footnote{在铁磁情形，要知道整个格子极化的净方向，只需知道一个自旋。在反铁磁情形，远处格点上自旋的方向当然还取决于它是否与原点的格点属于同一个子晶格。}

%FIG{189}: (a) 长程序蕴含短程序。(b) 短程序并不蕴含长程序。

我们前面已经指出，Bragg--Williams 方法甚至给不出任何短程序（short-range order），除非它是被假定的长程序所蕴含的：在居里温度以上，分布被假定得完全无规，每个自旋都与近邻的详细行为无关。Bethe 方法在 $T_c$ 以上会给出短程序，因此更为现实。两种方法都能相当好地描述长程有序态在充分建立之后的一般性质，但都不能用来证明这个态的确存在。
